\sectionheading{5.21 প্রোগ্রাম ডেভেলপমেন্ট (Program Development)}

একটি কার্যকর Program সাধারণত সরাসরি Code লেখা দিয়ে শুরু করা হয় না।
প্রথমে Problem নির্ধারণ ও বিশ্লেষণ করা হয়, তারপর Program Design,
Coding, Implementation, Documentation এবং Maintenance-এর মাধ্যমে
Program Development সম্পন্ন করা হয়।

%================================================
% Definition
%================================================
\booksubheading{Program Development কাকে বলে?}

\begin{definitionbox}

\textbf{\color{darkTeal}সংজ্ঞা}

কোনো Problem সমাধানের উদ্দেশ্যে Problem নির্ধারণ থেকে শুরু করে
Program Design, Coding, Implementation, Documentation এবং
Maintenance পর্যন্ত সম্পূর্ণ প্রক্রিয়াকে
\textbf{Program Development} বলে।

\end{definitionbox}

\textbf{সহজভাবে মনে রাখবে—}

\begin{center}

\textbf{
Problem বুঝি
\(\longrightarrow\)
Design করি
\(\longrightarrow\)
Code লিখি
\(\longrightarrow\)
Test করি
\(\longrightarrow\)
সংরক্ষণ ও উন্নয়ন করি
}

\end{center}

%================================================
% Main Steps
%================================================
\booksubheading{Program Development-এর প্রধান ধাপ}

Program তৈরির প্রধান ধাপগুলো হলো—

\begin{center}

\textbf{
Problem Identification
\(\rightarrow\)
Problem Analysis
\(\rightarrow\)
Program Design
}

\vspace{0.6mm}

\textbf{
\(\rightarrow\)
Development/Coding
\(\rightarrow\)
Implementation
\(\rightarrow\)
Documentation
\(\rightarrow\)
Maintenance
}

\end{center}

\begin{tabularx}{\linewidth}{|>{\bfseries}p{0.32\linewidth}|Y|}
\hline

\rowcolor{softBlue}
\multicolumn{1}{|c|}{\textbf{ধাপ}}
&
\multicolumn{1}{c|}{\textbf{মূল কাজ}}
\\
\hline

সমস্যা নির্দিষ্টকরণ
&
কোন Problem সমাধান করতে হবে তা নির্ধারণ করা।
\\
\hline

সমস্যা বিশ্লেষণ
&
Problem ও সম্ভাব্য Solution বিশ্লেষণ করা।
\\
\hline

Program Design
&
Algorithm ও Flowchart-এর সাহায্যে Program-এর পরিকল্পনা তৈরি করা।
\\
\hline

Development/Coding
&
Programming Language ব্যবহার করে Program Code লেখা।
\\
\hline

Program Implementation
&
Program পরীক্ষা, Error শনাক্ত ও প্রয়োজনীয় Correction করা।
\\
\hline

Program Documentation
&
Program সম্পর্কিত প্রয়োজনীয় লিখিত তথ্য সংরক্ষণ করা।
\\
\hline

Program Maintenance
&
প্রয়োজন অনুযায়ী Program পরিবর্তন, সংশোধন ও উন্নয়ন করা।
\\
\hline

\end{tabularx}

%================================================
% 1 Problem Identification
%================================================
\booksubheading{১. সমস্যা নির্দিষ্টকরণ (Problem Identification)}

Program Development-এর শুরুতে
\textbf{কোন Problem সমাধান করতে হবে তা সঠিকভাবে নির্ধারণ}
করা হয়।

এ পর্যায়ে Problem-এর উদ্দেশ্য ও কী ধরনের Result প্রয়োজন
সে সম্পর্কে পরিষ্কার ধারণা নেওয়া হয়।

\begin{center}

\textbf{
প্রথম প্রশ্ন:
``কোন Problem-এর সমাধান করতে হবে?''
}

\end{center}

%================================================
% 2 Problem Analysis
%================================================
\booksubheading{২. সমস্যা বিশ্লেষণ (Problem Analysis)}

Problem নির্ধারণের পর সেটি বিশ্লেষণ করা হয় এবং কীভাবে
Solution তৈরি করা সম্ভব তা বিবেচনা করা হয়।

একাধিক Solution সম্ভব হলে উপযুক্ত Solution নির্বাচন করা হয়।

Problem Analysis ও IPO Model আগের Topic-এ বিস্তারিত আলোচনা
করা হয়েছে; তাই এখানে পুনরাবৃত্তি করা হলো না।

%================================================
% 3 Program Design
%================================================
\booksubheading{৩. Program Design}

Problem Analysis-এর পর Program-এর সম্পূর্ণ পরিকল্পনা তৈরি করা হয়।

এই পর্যায়ে সাধারণত—

\begin{itemize}

\item Algorithm তৈরি করা হয়।

\item Flowchart তৈরি করা হয়।

\item Program-এর Logical Structure নির্ধারণ করা হয়।

\end{itemize}

\begin{rememberbox}

\textbf{\color{deepOrange}MCQ-এর জন্য গুরুত্বপূর্ণ}

Program তৈরির ক্ষেত্রে
\textbf{Program Design একটি অত্যন্ত গুরুত্বপূর্ণ ধাপ}।

Program-এর Solution কীভাবে বাস্তবায়িত হবে, তার পরিকল্পনা
এই ধাপেই তৈরি করা হয়।

\end{rememberbox}

%================================================
% 4 Coding
%================================================
\booksubheading{৪. Program Development বা Coding}

Program Design সম্পন্ন হওয়ার পর নির্বাচিত Programming Language
ব্যবহার করে Program-এর Instruction লেখা হয়।

Programming Language-এ Program লেখার এই কাজকে
\textbf{Coding} বলা হয়।

\begin{center}

\textbf{
Algorithm / Flowchart
\(\longrightarrow\)
Programming Language
\(\longrightarrow\)
Source Code
}

\end{center}

Coding-এর সময় Program-এর Logic ও Programming Language-এর
Syntax সঠিকভাবে ব্যবহার করতে হয়।

%================================================
% 5 Implementation
%================================================
\booksubheading{৫. Program Implementation}

Program লেখা শেষ হওয়ার পর Program-টি পরীক্ষা করা হয়।

Error পাওয়া গেলে প্রয়োজনীয় Correction করা হয় এবং Program
সঠিকভাবে কাজ করছে কি না যাচাই করা হয়।

\begin{center}
\begin{adjustbox}{max width=0.98\linewidth}
\fcolorbox{darkBlue}{softBlue}{
    \parbox[c][8mm][c]{0.15\linewidth}{
        \centering\textbf{Code}
    }
}
\(\longrightarrow\)
\fcolorbox{mainTeal}{softGreen}{
    \parbox[c][8mm][c]{0.15\linewidth}{
        \centering\textbf{Test}
    }
}
\(\longrightarrow\)
\fcolorbox{darkBlue}{softBlue}{
    \parbox[c][8mm][c]{0.17\linewidth}{
        \centering\textbf{Error\\Check}
    }
}
\(\longrightarrow\)
\fcolorbox{mainTeal}{softGreen}{
    \parbox[c][8mm][c]{0.18\linewidth}{
        \centering\textbf{Correction}
    }
}
\(\longrightarrow\)
\fcolorbox{darkBlue}{softBlue}{
    \parbox[c][8mm][c]{0.15\linewidth}{
        \centering\textbf{Retest}
    }
}
\end{adjustbox}
\end{center}

Testing ও Debugging এই পর্যায়ে গুরুত্বপূর্ণ ভূমিকা পালন করে।

%================================================
% Run Program
%================================================
\booksubheading{Run Program কী?}

Program-এর Error সংশোধন করার পর Program সঠিকভাবে কাজ করলে
Program-টি চালিয়ে Result পাওয়া যায়।

Program চালানোর প্রক্রিয়াকে সাধারণভাবে
\textbf{Run Program} বলা হয়।

\begin{center}

\textbf{
Error Correction
\(\longrightarrow\)
Run Program
\(\longrightarrow\)
Result
}

\end{center}

%================================================
% 6 Documentation
%================================================
\booksubheading{৬. Program Documentation}

\begin{definitionbox}

Program ভবিষ্যতে ব্যবহার, পরিবর্তন বা Update করার সুবিধার জন্য
Program সম্পর্কিত প্রয়োজনীয় তথ্য
\textbf{লিখিতভাবে সংরক্ষণ করার প্রক্রিয়াকে}
\textbf{Program Documentation} বলা হয়।

\end{definitionbox}

Documentation-এ প্রয়োজন অনুযায়ী থাকতে পারে—

\begin{itemize}

\item Program-এর উদ্দেশ্য
\item Program-এর Logic বা কাজের বিবরণ
\item প্রয়োজনীয় Input ও Output সম্পর্কিত তথ্য
\item ব্যবহারের নির্দেশনা
\item গুরুত্বপূর্ণ পরিবর্তনের তথ্য

\end{itemize}

ভালো Documentation থাকলে ভবিষ্যতে Program বোঝা,
ব্যবহার ও পরিবর্তন করা সহজ হয়।

%================================================
% 7 Maintenance
%================================================
\booksubheading{৭. Program Maintenance}

Program তৈরি ও ব্যবহার শুরু হওয়ার পরও প্রয়োজন অনুযায়ী
Program-এর বিভিন্ন পরিবর্তন করতে হতে পারে।

যেমন—

\begin{itemize}

\item Error Correction
\item Modification
\item পরিবর্ধন
\item প্রয়োজন অনুযায়ী Update বা উন্নয়ন

\end{itemize}

এই কাজগুলোকে \textbf{Program Maintenance} বলা হয়।

%================================================
% Documentation vs Maintenance
%================================================
\booksubheading{Documentation ও Maintenance-এর পার্থক্য}

\begin{tabularx}{\linewidth}{|>{\bfseries}p{0.24\linewidth}|Y|Y|}
\hline

\rowcolor{softBlue}
\multicolumn{1}{|c|}{\textbf{বিষয়}}
&
\multicolumn{1}{c|}{\textbf{Documentation}}
&
\multicolumn{1}{c|}{\textbf{Maintenance}}
\\
\hline

মূল কাজ
&
Program-এর তথ্য লিখিতভাবে সংরক্ষণ
&
Program পরিবর্তন ও উন্নয়ন
\\
\hline

উদ্দেশ্য
&
ভবিষ্যতে Program বোঝা ও ব্যবহার সহজ করা
&
Program কার্যকর ও প্রয়োজনোপযোগী রাখা
\\
\hline

উদাহরণ
&
Logic, ব্যবহারবিধি, পরিবর্তনের তথ্য
&
Error Correction, Modification, Update
\\
\hline

\end{tabularx}

%================================================
% Full Flow
%================================================
\booksubheading{Program Development-এর পূর্ণ Flow}

\begin{center}

\fcolorbox{darkBlue}{softBlue}{
    \parbox[c][8mm][c]{0.26\linewidth}{
        \centering\textbf{Problem Identification}
    }
}

\(\downarrow\)

\fcolorbox{mainTeal}{softGreen}{
    \parbox[c][8mm][c]{0.26\linewidth}{
        \centering\textbf{Problem Analysis}
    }
}

\(\downarrow\)

\fcolorbox{darkBlue}{softBlue}{
    \parbox[c][8mm][c]{0.26\linewidth}{
        \centering\textbf{Program Design}
    }
}

\(\downarrow\)

\fcolorbox{mainTeal}{softGreen}{
    \parbox[c][8mm][c]{0.26\linewidth}{
        \centering\textbf{Development / Coding}
    }
}

\(\downarrow\)

\fcolorbox{darkBlue}{softBlue}{
    \parbox[c][8mm][c]{0.26\linewidth}{
        \centering\textbf{Implementation}
    }
}

\(\downarrow\)

\fcolorbox{mainTeal}{softGreen}{
    \parbox[c][8mm][c]{0.26\linewidth}{
        \centering\textbf{Documentation}
    }
}

\(\downarrow\)

\fcolorbox{darkBlue}{softBlue}{
    \parbox[c][8mm][c]{0.26\linewidth}{
        \centering\textbf{Maintenance}
    }
}

\end{center}

%================================================
% Quick Comparison
%================================================
\booksubheading{ধাপগুলো দ্রুত মনে রাখো}

\begin{tabularx}{\linewidth}{|>{\bfseries}p{0.33\linewidth}|Y|}
\hline

\rowcolor{softBlue}
\multicolumn{1}{|c|}{\textbf{ধাপ}}
&
\multicolumn{1}{c|}{\textbf{সহজ প্রশ্ন}}
\\
\hline

Problem Identification
&
কোন Problem সমাধান করব?
\\
\hline

Problem Analysis
&
Problem কীভাবে সমাধান করা যাবে?
\\
\hline

Program Design
&
Solution-এর ধাপ কী হবে?
\\
\hline

Coding
&
Solution-কে কীভাবে Code-এ লিখব?
\\
\hline

Implementation
&
Program ঠিকভাবে কাজ করছে কি?
\\
\hline

Documentation
&
Program-এর প্রয়োজনীয় তথ্য কীভাবে সংরক্ষণ করব?
\\
\hline

Maintenance
&
ভবিষ্যতে কীভাবে Program সংশোধন বা উন্নয়ন করব?
\\
\hline

\end{tabularx}

%================================================
% Final Reminder
%================================================
\begin{rememberbox}

\textbf{\color{deepOrange}মনে রাখবে}

\begin{itemize}

\item Program Development হলো Program তৈরির
\textbf{সম্পূর্ণ প্রক্রিয়া}।

\item প্রথম ধাপ হলো
\textbf{Problem Identification}।

\item Problem Identification-এর পরে
\textbf{Problem Analysis} করা হয়।

\item Program Design-এ Algorithm ও Flowchart ব্যবহার করে
Solution-এর পরিকল্পনা করা হয়।

\item Programming Language-এ Program লেখার কাজকে
\textbf{Coding} বলা হয়।

\item Implementation পর্যায়ে Program
\textbf{Test ও Correct} করা হয়।

\item \textbf{Documentation} =
Program সম্পর্কিত প্রয়োজনীয় লিখিত তথ্য সংরক্ষণ।

\item \textbf{Maintenance} =
Program-এর প্রয়োজনীয় পরিবর্তন, সংশোধন ও উন্নয়ন।

\item দ্রুত ক্রম—

\begin{center}

\textbf{
Identify
\(\rightarrow\)
Analyze
\(\rightarrow\)
Design
\(\rightarrow\)
Code
}

\vspace{0.5mm}

\textbf{
\(\rightarrow\)
Implement
\(\rightarrow\)
Document
\(\rightarrow\)
Maintain
}

\end{center}

\end{itemize}

\end{rememberbox}

\small

%================================================
% 5.22 Dry Run and Trace Table
%================================================
